\documentclass[uplatex,dvipdfmx,a4paper,11pt]{jsarticle}
\input{../../../templates/kaisetsu.tex}
\kaisetsuprint{量子力学1　第1章「量子論の誕生（黒体輻射・光の粒子性・物質波・Bohr の原子論）」}{2026-11-24}
%% kaisetsu-moto	problems\量子力学1\shou01\量子力学1-1.tex	kiroku\kaisetsu\量子力学1\量子力学1-1_2026-11-24_komento.tex
\begin{document}
\matome
\begin{saitenichiran}
  \hyougyou{\hyperlink{bt-m-119-bohr}{7}}{L2・計算}{△△△△}{C}{ヒント1を見た / (1)計算 (2)計算 (3)計算 (4)計算}
\end{saitenichiran}
\begin{daiji}
  \item 円運動する電子の力学的エネルギーは位置エネルギーの半分、$E_{n}=-e^{2}/(8\pi\varepsilon_{0}r_{n})$ である。この係数 $1/2$ を落とすと、$E_{1}$ は2倍、波長は半分になり、(2)〜(4) の数値がすべてずれる（問題7 (1)〜(4)）。
  \item 束縛状態のエネルギーは負である。$E_{1}=\qty{-13.6}{\electronvolt}$（水素のイオン化エネルギー）、$\lambda_{32}=\qty{656}{\nano\metre}$（赤い Hα 線）のようなよく知られた値で、数値の答えを検算する（問題7 (2)〜(4)）。
\end{daiji}
\begin{maemo}
  \item 同じ問題 (1) の $E_{n}$ の計算で、2026-10-29 にも途中で係数を書き違えた（そのときは次の行で直した）。今回は直らずに最後まで残った。$E_{n}=K+U=U/2$ を1行書いてから $r_{n}$ を代入する。
\end{maemo}

\mondaihen
\level{2}
\mondaibangou{7}
\begin{mondai}{119-bohr}{計算}{Bohr の原子論（量子条件と振動数条件）による水素原子の軌道半径・エネルギー準位と，水素のスペクトルの波長}
\begin{teisuu}
  プランク定数 & $h$ & \qty{6.626e-34}{\joule\second}\\
  光速 & $c$ & \qty{2.998e8}{\metre\per\second}\\
  電気素量 & $e$ & \qty{1.602e-19}{\coulomb}\\
  電子の質量 & $\melec$ & \qty{9.109e-31}{\kilogram}\\
  陽子の質量 & $\mprot$ & \qty{1.673e-27}{\kilogram}\\
  ボルツマン定数 & $\kB$ & \qty{1.381e-23}{\joule\per\kelvin}\\
  電子の静止エネルギー & $\melec c^{2}$ & \qty{0.511}{\mega\electronvolt}\\
  真空の誘電率 & $\varepsilon_{0}$ & \qty{8.854e-12}{\farad\per\metre}\\
\end{teisuu}
水素原子を，静止した陽子（電荷 $e$）のまわりを電子（質量 $\melec$，電荷 $-e$）が半径 $r$ の等速円運動をするものとして，要点整理の Bohr の原子論で考える。数値は定数表の値と $\hbar=h/(2\pi)$ を使い，$\qty{1}{\electronvolt}=\qty{1.602e-19}{\joule}$ とする。
  \begin{shoumon}
    \item クーロン力が向心力になる条件と量子条件から，許される軌道の半径 $r_{n}$ と，電子の力学的エネルギー（運動エネルギーとクーロン力の位置エネルギーの和）$E_{n}$ を，$n$，$\melec$，$e$，$\varepsilon_{0}$，$\hbar$ で表せ。
    \item $r_{1}$（Bohr 半径）を \unit{\nano\metre} 単位で，$E_{1}$ を \unit{\electronvolt} 単位で，それぞれ有効数字3桁で求めよ。
    \item 電子が $n=3$ の軌道から $n=2$ の軌道に移るときに放出される光の波長を，有効数字3桁で求めよ。
    \item $n\ge3$ の軌道から $n=2$ の軌道に移るときに出る光の波長には下限がある。その値（$n\to\infty$ の極限）を有効数字3桁で求めよ。
  \end{shoumon}

\begin{kotaeran}
  \machigai{1}{△}{計算} 運動方程式と量子条件から $r_{n}=4\pi\varepsilon_{0}n^{2}\hbar^{2}/(\melec e^{2})=n^{2}\varepsilon_{0}h^{2}/(\pi\melec e^{2})$ は正しい。$E_{n}=e^{2}/(8\pi\varepsilon_{0}r_{n})-e^{2}/(4\pi\varepsilon_{0}r_{n})$ までも正しいが、次の行で $-\bigl(e^{2}/(4\pi\varepsilon_{0})\bigr)\cdot\pi\melec e^{2}/(n^{2}\varepsilon_{0}h^{2})$ と係数 $1/2$ が落ち、$E_{n}=-\melec e^{4}/(4\varepsilon_{0}^{2}h^{2}n^{2})$ とした。$1/2-1=-1/2$ なので $E_{n}=-e^{2}/(8\pi\varepsilon_{0}r_{n})=-\melec e^{4}/(8\varepsilon_{0}^{2}h^{2}n^{2})$ が正しい。
  \machigai{2}{△}{計算} $r_{1}=\qty{0.0529}{\nano\metre}$ は正しい。$E_{1}$ を \qty{27.2}{\electronvolt} とした。(1) の係数の誤りで大きさが2倍になり、さらに式の負号を落として正の値にしている。正しくは $E_{1}=\qty{-13.6}{\electronvolt}$。
  \machigai{3}{△}{計算} 振動数条件 $\Delta E=hc/\lambda$ を使う方針は正しい。(2) の $E_{1}$ を正として $E_{3}-E_{2}=\frac{5}{36}E_{1}=\qty{3.78}{\electronvolt}$、$\lambda=\qty{328}{\nano\metre}$ とした（(1) の2倍を引き継いで、正しい値の半分）。$E_{1}<0$ なので $E_{3}-E_{2}=-\frac{5}{36}E_{1}=\frac{5}{36}\abs{E_{1}}=\qty{1.89}{\electronvolt}$、$\lambda=\qty{656}{\nano\metre}$。
  \machigai{4}{△}{計算} $0-E_{2}=\frac{1}{4}E_{1}$ とし $\lambda=\qty{182}{\nano\metre}$ とした（(3) と同じく2倍の $E_{1}$ を引き継いだ）。正しくは $0-E_{2}=\abs{E_{1}}/4$、$\lambda=\qty{365}{\nano\metre}$。また、$n$ が大きいほど $E_{n}-E_{2}$ が大きく波長が短いので下限が $n\to\infty$ の極限になる、という根拠が書かれていない（解答の論点）。
\end{kotaeran}
\begin{rikai}
  \item 運動方程式 $\melec v^{2}/r=e^{2}/(4\pi\varepsilon_{0}r^{2})$ の両辺に $r/2$ をかけると、運動エネルギー $K=e^{2}/(8\pi\varepsilon_{0}r)=-U/2$。よって $E=K+U=U/2=-e^{2}/(8\pi\varepsilon_{0}r)$（ビリアル定理の形）。「$E$ は $U$ の半分」と言葉で押さえておけば、係数の書き落としに気づける。
  \item 束縛状態の $E_{n}=E_{1}/n^{2}$ は $E_{1}<0$ のまま使い、遷移のエネルギーは $E_{3}-E_{2}=\abs{E_{1}}(1/4-1/9)$ のように絶対値で書くと、符号を取り違えない。
  \item 検算：\qty{27.2}{\electronvolt} は \qty{13.6}{\electronvolt} のちょうど2倍、\qty{328}{\nano\metre} は \qty{656}{\nano\metre} のちょうど半分である。答えが既知の値の2倍・半分になったら、まず係数 $1/2$ を疑う。
\end{rikai}
\end{mondai}

\kaitouhen
\level{2}
\kaisetsubangou{7}
\begin{kaisetsu}{119-bohr}
\houshin 円運動の運動方程式（クーロン力が向心力）と量子条件 $\melec vr=n\hbar$ を連立して $r$ と $v$ を決め，力学的エネルギーを求める。遷移で出る光の波長は，振動数条件 $h\nu=E_{n}-E_{n'}$ と $\lambda=c/\nu$ から求める。
\keisan
\shoumonkai{1} 運動方程式と量子条件は
\[ \frac{\melec v^{2}}{r}=\frac{e^{2}}{4\pi\varepsilon_{0}r^{2}},\qquad \melec vr=n\hbar \]
である。第2式の $v=n\hbar/(\melec r)$ を第1式に代入すると
\[ \frac{n^{2}\hbar^{2}}{\melec r^{3}}=\frac{e^{2}}{4\pi\varepsilon_{0}r^{2}}\quad\text{より}\quad r_{n}=\frac{4\pi\varepsilon_{0}\hbar^{2}}{\melec e^{2}}\,n^{2} \]
運動方程式の両辺に $r/2$ をかけると，運動エネルギーは $\melec v^{2}/2=e^{2}/(8\pi\varepsilon_{0}r)$ で，位置エネルギー $-e^{2}/(4\pi\varepsilon_{0}r)$ の $-1/2$ 倍である。よって
\[ E_{n}=\frac{e^{2}}{8\pi\varepsilon_{0}r_{n}}-\frac{e^{2}}{4\pi\varepsilon_{0}r_{n}}=-\frac{e^{2}}{8\pi\varepsilon_{0}r_{n}}=-\frac{\melec e^{4}}{2(4\pi\varepsilon_{0})^{2}\hbar^{2}}\,\frac{1}{n^{2}} \]
\shoumonkai{2} $\hbar=h/(2\pi)=\qty{1.0546e-34}{\joule\second}$ を使うと
\begin{align*}
  r_{1}&=\frac{4\pi\times\num{8.854e-12}\times(\num{1.0546e-34})^{2}}{\num{9.109e-31}\times(\num{1.602e-19})^{2}}\,\unit{\metre}=\qty{5.29e-11}{\metre}=\qty{0.0529}{\nano\metre},\\
  E_{1}&=-\frac{e^{2}}{8\pi\varepsilon_{0}r_{1}}=-\frac{(\num{1.602e-19})^{2}}{8\pi\times\num{8.854e-12}\times\num{5.293e-11}}\,\unit{\joule}=\qty{-2.179e-18}{\joule}=\qty{-13.6}{\electronvolt}.
\end{align*}
\shoumonkai{3} $E_{n}=E_{1}/n^{2}$ なので，放出される光子のエネルギーは
\[ h\nu=E_{3}-E_{2}=\abs{E_{1}}\left(\frac{1}{4}-\frac{1}{9}\right)=\frac{5}{36}\abs{E_{1}} \]
で，波長は
\[ \lambda_{32}=\frac{hc}{E_{3}-E_{2}}=\frac{36hc}{5\abs{E_{1}}}=\frac{36\times\num{6.626e-34}\times\num{2.998e8}}{5\times\num{2.179e-18}}\,\unit{\metre}=\qty{6.56e-7}{\metre}=\qty{656}{\nano\metre} \]
\shoumonkai{4} $n$ が大きいほど $E_{n}-E_{2}$ は大きく，波長は短い。$n\to\infty$ で $E_{n}\to0$ なので，波長の下限は
\[ \lambda_{\infty}=\frac{hc}{0-E_{2}}=\frac{4hc}{\abs{E_{1}}}=\frac{4\times\num{6.626e-34}\times\num{2.998e8}}{\num{2.179e-18}}\,\unit{\metre}=\qty{365}{\nano\metre} \]
\zu[水素原子のエネルギー準位と，$n=2$ への遷移（矢印）]{\begin{tikzpicture}[x=1cm,y=0.22cm]
  \draw[jiku] (0,-14.8) -- (0,2.2) node[above]{$E$／\unit{\electronvolt}};
  \foreach \E/\n in {-13.6/1,-3.40/2,-1.51/3}{\draw[junni] (0,\E) -- (4.5,\E) node[right]{$n=\n$};}
  \draw[junni] (0,-0.85) -- (4.5,-0.85);
  \draw[junni] (0,0) -- (4.5,0) node[right]{$n=\infty$};
  \foreach \E in {-13.6,-3.40,-1.51}{\node[left,font=\small] at (0,\E) {$\E$};}
  \node[left,font=\small] at (0,0) {$0$};
  \draw[yajirushi] (1.5,-1.51) -- (1.5,-3.40) node[midway,right,font=\small]{\qty{656}{\nano\metre}};
  \draw[yajirushi] (3.2,0) -- (3.2,-3.40) node[midway,right,font=\small]{\qty{365}{\nano\metre}};
\end{tikzpicture}}
\kotae{\shoumonkai{1} $r_{n}=\frac{4\pi\varepsilon_{0}\hbar^{2}}{\melec e^{2}}n^{2}$，$E_{n}=-\frac{\melec e^{4}}{2(4\pi\varepsilon_{0})^{2}\hbar^{2}n^{2}}$ \shoumonkai{2} $r_{1}=\qty{0.0529}{\nano\metre}$，$E_{1}=\qty{-13.6}{\electronvolt}$ \shoumonkai{3} \qty{656}{\nano\metre} \shoumonkai{4} \qty{365}{\nano\metre}}
\begin{ronten}
  \item \shoumonkai{1} 力学的エネルギーを運動エネルギーと位置エネルギー（無限遠を基準）の和とし，運動方程式から運動エネルギーが位置エネルギーの $-1/2$ 倍になることを使ったと書く。
  \item \shoumonkai{4} $n$ が大きいほど $E_{n}-E_{2}$ が大きく波長が短いこと（単調性）から，下限が $n\to\infty$ の極限で与えられることを述べる。
\end{ronten}
\kenzan 次元：$\varepsilon_{0}$ は（電荷）$^{2}$／（エネルギー×長さ）の次元をもつので，$\varepsilon_{0}\hbar^{2}/(\melec e^{2})$ は（エネルギー×時間）$^{2}$／（質量×エネルギー×長さ）になり，エネルギー＝質量×（長さ／時間）$^{2}$ を使うと長さになる。$E_{n}$ は $e^{2}/(\varepsilon_{0}r)$ の形でエネルギーである。$\lambda_{32}=\qty{656}{\nano\metre}$ は赤色の可視光（水素の Hα 線）で，$\lambda_{\infty}$ はそれより短い。また $v_{1}=\hbar/(\melec r_{1})=\qty{2.19e6}{\metre\per\second}$ は光速の約 $1/137$ で，非相対論的に扱ってよい。
\tsumazuki 力学的エネルギーは運動エネルギーと位置エネルギーの和で，運動エネルギーは位置エネルギーの大きさの半分である。位置エネルギーだけを $E_{n}$ としない。
\end{kaisetsu}

\end{document}
